%source mathadata
On considère deux nuages de points représentés dans le repère ci-contre : les points du nuage des 2 et ceux du nuage des 7. Les points de référence sont $C(2\,;\,4)$ pour le chiffre 2 et $R(5\,;\,2)$ pour le chiffre 7.

		\begin{multicols}{2}

			\begin{MADenumerate}

				\item Calculer les distances $TC$, $TR$, $UC$ et $UR$ pour les points $T(4\,;\,3)$ et $U(3\,;\,2{,}5)$. Quel chiffre l'algorithme estime-t-il pour chacun ?

				\item On cherche l'équation de la droite qui sépare les deux zones d'estimation.

				Soit $M(x;y)$ un point sur cette droite.
				\begin{enumerate}
					\item Que peut-on dire des distances $MC$ et $MR$ ?
					\item Justifier que :
					      $$(x-2)^2+(y-4)^2=(x-5)^2+(y-2)^2$$
					\item Déduire l'équation de la droite sous la forme $y=mx+p$.
				\end{enumerate}

				\begin{exemple}[Aide]
					Vous devriez :
					\begin{enumerate}
						\item Développer l'équation de la question b.
						\item Simplifier.
						\item Mettre l'équation sous la forme $y=mx+p$.
					\end{enumerate}
				\end{exemple}

				\item Calculer les coordonnées de deux points de la droite, par exemple pour $x=0$ et $x=2$.

				\item Tracer la droite dans le repère ci-contre. Puis, colorer les deux zones — une pour le chiffre 2, une pour le chiffre 7.

				\item Sans calculer de distances, quel chiffre l'algorithme estime-t-il pour le point $V(1\,;\,5)$ ? Justifier en utilisant les zones.

			\end{MADenumerate}

			\begin{figure}[H]
				\centering
				\begin{tikzpicture}[line cap=round,line join=round,>=triangle 45,x=1cm,y=1cm,scale=0.7]
					\begin{axis}[
							x=1cm,y=1cm,
							axis lines=middle,
							ymajorgrids=true,
							xmajorgrids=true,
							xmin=-1,
							xmax=8,
							ymin=-1,
							ymax=8,
							xtick={0,1,...,7},
							ytick={0,1,...,7},]
					\end{axis}
				\end{tikzpicture}
			\end{figure}

		\end{multicols}
